\section*{Capítulo \ref{ch:b2:vertebrate-development} --- Desenvolvimento embrionário dos vertebrados}

\begin{solution}{exo:b2:vertebrate-development:1}
\omterm{def:b2:vertebrate-development:egg}{Clivagem}: mitoses rápidas sem crescimento, que produzem a \omterm{def:b2:vertebrate-development:blastula}{blástula}.
\omterm{def:b2:vertebrate-development:blastula}{Blástula}: uma esfera oca ou um disco de células pequenas em torno da
\omterm{def:b2:vertebrate-development:blastula}{blastocela}. Gastrulação: os movimentos celulares que levam o mesoderma e
o endoderma para dentro, produzindo a gástrula de três camadas, com uma
cavidade intestinal e a notocorda. Neurulação: o dobramento do
ectoderma dorsal no \omterm{prop:b2:vertebrate-development:neurulation}{tubo neural}, com a segmentação dos somitos,
produzindo a nêurula com o plano corporal dos vertebrados.
\end{solution}

\begin{solution}{exo:b2:vertebrate-development:2}
Ectoderma: epiderme, encéfalo e medula espinal, nervos periféricos
(\omterm{prop:b2:vertebrate-development:neurulation}{crista neural}), cristalino e células pigmentares. Mesoderma:
notocorda, músculo, osso e cartilagem, rim, coração e sangue, derme.
Endoderma: o revestimento do intestino, fígado, pâncreas, pulmões,
tireoide.
\end{solution}

\begin{solution}{exo:b2:vertebrate-development:3}
Pouco \omterm{def:b2:vertebrate-development:egg}{vitelo}: \omterm{def:b2:vertebrate-development:egg}{clivagem} completa, quase igual, e gastrulação por
involução através de um blastóporo redondo numa esfera oca (rã:
completa, mas desigual, com as células vegetativas ricas em \omterm{def:b2:vertebrate-development:egg}{vitelo}
grandes e lentas). Muito \omterm{def:b2:vertebrate-development:egg}{vitelo} (galinha): \omterm{def:b2:vertebrate-development:egg}{clivagem} restrita a um disco
sobre o \omterm{def:b2:vertebrate-development:egg}{vitelo} e gastrulação através de uma fenda, a linha primitiva,
no disco achatado, com as camadas se espalhando sobre o \omterm{def:b2:vertebrate-development:egg}{vitelo} em vez
de envolverem uma cavidade.
\end{solution}

\begin{solution}{exo:b2:vertebrate-development:4}
A \omterm{def:b2:vertebrate-development:induction}{indução} é a mudança do destino de um tecido competente por um sinal
vindo de um tecido vizinho. As células vegetativas induzem o mesoderma
na zona marginal situada acima delas; o \omterm{def:b2:vertebrate-development:induction}{organizador} (mesoderma do lábio
dorsal) induz a placa neural no ectoderma dorsal; a vesícula óptica
induz o cristalino no ectoderma da cabeça.
\end{solution}

\begin{solution}{exo:b2:vertebrate-development:5}
$2^{12} = 4096$ células de raio $1/2^{4} = \qty{62.5}{\micro m}$;
superfície multiplicada por $2^{4} = 16$. A superfície adicional é a
membrana dos epitélios que a gastrulação dobrará em pele e intestino,
e a área através da qual a \omterm{def:b2:vertebrate-development:blastula}{blástula} faz trocas com o meio.
\end{solution}

\begin{solution}{exo:b2:vertebrate-development:6}
$10^{-4}\times 2^{k} = 0.4$: $2^{k} = 4000$, $k = 12$. Com quatro
vezes mais DNA, a razão começa em $4\times 10^{-4}$ e atinge $0.4$ em
$2^{k} = 1000$: a décima divisão, duas antes.
\end{solution}

\begin{solution}{exo:b2:vertebrate-development:7}
$k = \ln 2/600 = \qty{1.16e-3}{s^{-1}}$; $\lambda = \sqrt{2\times
10^{-12}/1.16\times 10^{-3}} = \qty{42}{\micro m}$. Limiares em
$x = \lambda\ln 2 = \qty{29}{\micro m}$ e $\lambda\ln 10 =
\qty{96}{\micro m}$: faixas de 29, 67 e o restante.
\end{solution}

\begin{solution}{exo:b2:vertebrate-development:8}
Dobrar a fonte desloca as duas fronteiras para fora em $\lambda\ln
2 = \qty{29}{\micro m}$: a primeira faixa dobra, a segunda conserva sua
largura. Dobrar $k$ divide $\lambda$ por $\sqrt 2$
(\qty{29}{\micro m}) e, como $c_0 = j/\sqrt{Dk}$, divide também $c_0$
por $\sqrt 2$: as fronteiras passam a $\lambda'\ln(c_0'/T) =
29\times(0.69 - 0.35) = \qty{10}{\micro m}$ e $29\times(2.30 -
0.35) = \qty{57}{\micro m}$ --- as duas faixas encolhem.
\end{solution}

\begin{solution}{exo:b2:vertebrate-development:9}
Um lábio dorsal de uma gástrula de tritão não pigmentada foi enxertado
no lado ventral de uma gástrula pigmentada; ele rolou para dentro e o
hospedeiro formou um segundo embrião --- \omterm{prop:b2:vertebrate-development:neurulation}{tubo neural}, somitos,
intestino --- unido ao primeiro. A diferença de pigmento mostrou que o
segundo eixo era feito de células do hospedeiro, de modo que o enxerto
as tinha induzido, em vez de construir ele mesmo o eixo. O resultado
excluiu a autodiferenciação do enxerto como explicação e mostrou que
uma pequena região leva o sinal que organiza o eixo inteiro.
\end{solution}

\begin{solution}{exo:b2:vertebrate-development:10}
(a) Pele do ventre: o ectoderma inicial é competente, mas ainda não foi
induzido, de modo que segue seu novo ambiente. (b) Tecido neural: na
gástrula tardia, ele já foi induzido e está determinado. (c) Placa
neural: o ectoderma ventral da gástrula tardia ainda é competente, e a
notocorda o induz. (d) Nenhuma estrutura dorsal: uma bola de tecido
ventral, pois nada induz o mesoderma a destinos dorsais nem o
ectoderma a destinos neurais.
\end{solution}

\begin{solution}{exo:b2:vertebrate-development:11}
A \qty{18}{\celsius}: $90 + 11\times 30 = \qty{420}{min}$, 7 horas.
A \qty{10}{\celsius}: \qty{17.5}{h}. A transição ocorre na décima
segunda divisão nos dois casos porque o que a desencadeia é a razão
entre o DNA e o citoplasma, que depende do número de divisões e não de
quanto tempo elas levaram: o relógio conta duplicações, não horas.
\end{solution}

\begin{solution}{exo:b2:vertebrate-development:12}
No desenvolvimento regulativo, as células adquirem seus destinos por
\omterm{def:b2:vertebrate-development:induction}{indução} a partir das vizinhas, de modo que o embrião consegue se
reconstruir após uma perda (gêmeos a partir de um ovo dividido,
regeneração de uma região enxertada); no desenvolvimento em mosaico, os
destinos são herdados com determinantes citoplasmáticos localizados,
de modo que cada blastômero forma apenas a sua parte e um embrião
dividido forma duas meias larvas. A diferença é de grau: o crescente
cinzento da rã é um determinante, e os embriões em mosaico usam
\omterm{def:b2:vertebrate-development:induction}{induções} mais tarde; todo embrião mistura herança e diálogo, e os
vertebrados se apoiam no diálogo por mais tempo.
\end{solution}

\begin{solution}{pb:b2:vertebrate-development:1}
\textbf{1.} Ovo $\tfrac43\pi(0.7)^{3} = \qty{1.44}{mm^3}$; núcleo
$\tfrac43\pi(0.03)^{3} = \qty{1.1e-4}{mm^3}$; razão $8\times
10^{-5}$.
\textbf{2.} $2^{k} = 4000$: $k = 12$, 4096 células.
\textbf{3.} $75 + 11\times 30 = \qty{405}{min}$, cerca de 6.75 horas.
\textbf{4.} Raio $0.7/2^{4} = \qty{44}{\micro m}$; razão $8\times
10^{-5}\times 4096 = 0.32$ --- a de uma célula comum: os núcleos
alcançaram o citoplasma.
\textbf{5.} $2^{4} = 16$.
\textbf{6.} Cerca de $4095\,c$ de DNA; o ovo contém histonas,
polimerases e nucleotídeos armazenados suficientes para doze divisões,
de modo que nenhuma transcrição é necessária.
\textbf{7.} Os genes zigóticos são ativados, as divisões desaceleram e
perdem a sincronia, as células se tornam móveis. Injetar DNA extra
antecipa a transição tantas divisões quanto o excesso de DNA, o que um
contador de divisões ou de tempo não conseguiria explicar.
\textbf{8.} Do lado oposto ao ponto de entrada do espermatozoide, no
crescente cinzento fixado pela rotação do córtex do ovo após a
fecundação.
\textbf{9.} Primeiro o mesoderma dorsal (placa pré-cordal, depois
notocorda), para o teto do arquêntero, na linha média; o mesoderma
lateral (somitos, placa lateral) ao lado dele; o endoderma reveste o
arquêntero; o ectoderma se espalha por fora por epibolia.
\textbf{10.} $\qty{2000}{\micro m}/\qty{600}{min} = \qty{3.3}{\micro
m/min}$: três vezes um fibroblasto --- uma lâmina coordenada se move
mais depressa que uma célula isolada.
\textbf{11.} Os movimentos celulares exigem produtos de genes
zigóticos (novas moléculas de adesão, motilidade) e ciclos mais lentos;
antes da transição, as células só conseguem se dividir.
\textbf{12.} Todo cordado tem uma notocorda em algum estágio, e os
vertebrados constroem a coluna vertebral em torno dela; o \omterm{prop:b2:vertebrate-development:neurulation}{tubo neural} é
induzido por ela e é a consequência, e não a origem, do plano.
\textbf{13.} Sem a contração das células em garrafa, nenhum blastóporo
se forma, nada rola para dentro e o embrião permanece uma esfera com as
camadas no lugar: sem intestino, sem notocorda, sem eixo.
\textbf{14.} $k = \ln 2/1200 = \qty{5.8e-4}{s^{-1}}$; $\lambda =
\sqrt{10^{-12}/5.8\times 10^{-4}} = \qty{42}{\micro m}$.
\textbf{15.} $x_1 = 42\ln 2.5 = \qty{38}{\micro m}$; $x_2 = 42\ln 12.5
= \qty{105}{\micro m}$.
\textbf{16.} Faixas de \qty{38}{\micro m} (2.5 células), \qty{67}{\micro
m} (4.5 células) e o restante.
\textbf{17.} $x_1 = 42\ln 1.25 = \qty{9}{\micro m}$; $x_2 = 42\ln 6.25
= \qty{77}{\micro m}$: as duas fronteiras se deslocam para dentro em
$\lambda\ln 2 =
\qty{29}{\micro m}$; a primeira faixa encolhe de 38 para 9, a faixa do
meio conserva sua largura, a faixa externa cresce.
\textbf{18.} Deslocada antes da leitura, ela adota o destino A --- lê
a concentração do lugar onde agora está; deslocada depois, conserva o
destino B --- já estava determinada.
\textbf{19.} $x_T = \lambda\ln(c_0/T)$: uma mudança da fonte por um
fator $f$ desloca todas as fronteiras em $\lambda\ln f$, de modo que um
erro de um fator dois desloca o padrão em apenas $0.7\lambda$, uma
fração de um campo que se estende por vários $\lambda$ --- o padrão é
quase insensível à quantidade exata de morfógeno.
\textbf{20.} Lábio dorsal no lado ventral: um segundo eixo. Zona
marginal ventral no lado dorsal: ela é reespecificada pelo ambiente e
forma tecido dorsal; nada a mais.
\textbf{21.} Numa nêurula, o ectoderma já não é competente: o enxerto
forma notocorda, mas não induz um segundo \omterm{prop:b2:vertebrate-development:neurulation}{tubo neural}.
\textbf{22.} Separado no plano do crescente: dois embriões completos;
perpendicularmente: um embrião e uma bola de tecido ventral.
\textbf{23.} Sozinha: uma bola epidérmica. Com células vegetativas:
mesoderma --- músculo, notocorda --- induzido na calota.
\textbf{24.} Para provar que o segundo eixo era induzido em células do
hospedeiro, e não construído pelo enxerto; se o enxerto fosse
indistinguível, a autodiferenciação não poderia ter sido excluída. Um
enxerto de outra espécie funciona (eles usaram duas espécies de
tritão), e foi isso que tornou as células distinguíveis.
\textbf{25.} Décima segunda \omterm{def:b2:vertebrate-development:egg}{clivagem}, 4096 células, por volta de 6.75
horas; faixas de 38 e \qty{67}{\micro m}, e tudo o que fica além de
\qty{105}{\micro m}.
\end{solution}
